Na frente de aprendizado de máquina, um ARIMA$(0,1,1)$ por canal de frequência
previu o RSSI de um nó sensor real um passo à frente com RMSE de 1,091~dB e MAE
de 0,813~dB no teste, contra 1,319 e 0,902~dB do \textit{random walk} por canal
(Seção~\ref{subsec:ml-res-frequencia}). A maior parte do erro que sobrava no
modelo da série inteira veio da troca de canal a cada quadro, já que cada canal
teve um nível próprio de RSSI. Com 14 dias de treino, cerca de 500 quadros por
canal, o modelo atendeu ao critério em todos os períodos avaliados, nos dois
\textit{gateways} estudados. Com cerca de 1\% de alarmes falsos, a regra baseada
no erro de previsão detectou, em 300 posições do teste, todos os desvios
sintéticos de 5~dB ou mais que duraram cerca de 1~h, e de 91\% a 92\% dos
desvios pontuais de 5~dB (Seção~\ref{subsec:ml-res-deteccao}). Em desvios
pontuais, o modelo detectou mais que um detector baseado no \textit{random walk}
por canal, mesmo quando este tinha quase o dobro de alarmes falsos
(Seção~\ref{subsec:ml-res-comparacao}). Esses resultados medem a sensibilidade
a desvios sintéticos, e não a detecção de ataques reais.
